\documentclass[11pt,a4paper]{article}
\usepackage[utf8]{inputenc}
\usepackage[spanish,es-tabla]{babel}
\usepackage[margin=2.2cm,top=2.5cm,bottom=2.5cm]{geometry}
\usepackage{graphicx}
\usepackage{xcolor}
\usepackage{titlesec}
\usepackage{booktabs}
\usepackage{tabularx}
\usepackage{array}
\usepackage{enumitem}
\usepackage{fancyhdr}
\usepackage{tcolorbox}
\usepackage{listings}
\usepackage{hyperref}

% Colores institucionales del equipo Los Gantt-sters
\definecolor{gblue}{RGB}{27, 79, 114}     % Azul corporativo
\definecolor{gnavy}{RGB}{21, 67, 96}      % Azul oscuro
\definecolor{ggreen}{RGB}{39, 174, 96}    % Verde éxito
\definecolor{ggray}{RGB}{245, 247, 250}   % Gris fondo
\definecolor{gborder}{RGB}{208, 211, 212} % Gris borde
\definecolor{gorange}{RGB}{211, 84, 0}    % Naranja advertencia

\hypersetup{
    colorlinks=true,
    linkcolor=gblue,
    urlcolor=gblue,
    citecolor=gblue
}

% Cabeceras y pies
\pagestyle{fancy}
\fancyhf{}
\fancyhead[L]{\small\color{gray} Los Gantt-sters · Agensible (TuAgendita)}
\fancyhead[R]{\small\color{gray} Estandarización Técnica Sprint 1}
\fancyfoot[L]{\small\color{gray} DGP \& MDA --- UGR · Curso 2026/2027}
\fancyfoot[R]{\small\color{gray} Página \thepage}
\renewcommand{\headrulewidth}{0.4pt}
\renewcommand{\footrulewidth}{0.4pt}

% Formato de títulos
\titleformat{\section}{\Large\bfseries\color{gblue}}{\thesection.}{0.8em}{}[\titlerule]
\titleformat{\subsection}{\large\bfseries\color{gnavy}}{\thesubsection.}{0.6em}{}
\titleformat{\subsubsection}{\normalsize\bfseries\color{gblue}}{\thesubsubsection.}{0.5em}{}

\begin{document}

% PORTADA OFICIAL
\begin{titlepage}
    \centering
    \vspace*{0.5cm}
    \includegraphics[width=0.36\textwidth]{logo.png}\\[1cm]
    
    {\Huge \bfseries \color{gblue} AGENSIBLE (TuAgendita)}\\[0.3cm]
    {\large \bfseries \color{gray} Agenda Personalizada y Accesible para Estudiantes con Discapacidad Intelectual}\\[0.2cm]
    {\normalsize CEE Fundación Purísima Concepción de Granada}\\[1.2cm]
    
    \begin{tcolorbox}[colback=gblue,colframe=gnavy,arc=4mm,width=0.95\textwidth,center]
        \color{white}
        \centering
        \vspace{0.2cm}
        {\LARGE \bfseries ESTÁNDAR TÉCNICO OFICIAL Y STACK COMÚN}\\[0.2cm]
        {\large Revisión de Tareas, Directrices de Desarrollo y Homogeneización del Sprint 1}\\[0.2cm]
    \end{tcolorbox}
    
    \vspace{1.5cm}
    
    \begin{minipage}{0.85\textwidth}
        \centering
        \textbf{Asignaturas:} Dirección y Gestión de Proyectos (DGP) \&\\ Metodologías de Desarrollo Ágil (MDA)\\[0.2cm]
        \textbf{Grado:} Grado en Ingeniería Informática --- Universidad de Granada\\[0.2cm]
        \textbf{Subgrupo 2:} \textit{Los Gantt-sters}\\[0.5cm]
        \textbf{Coordinador del Proyecto:} Juan Luis Jiménez Urbano\\
        \textbf{Fecha:} Noviembre 2026
    \end{minipage}
    
    \vfill
    {\footnotesize Documento de uso interno para la coordinación técnica y estandarización del equipo de desarrollo.}
\end{titlepage}

\newpage

% ÍNDICE
\tableofcontents
\vspace{1cm}
\hrule
\vspace{1cm}

\section{Propósito y Justificación de la Estandarización}

Durante las primeras sesiones de desarrollo convivieron en el repositorio dos enfoques técnicos: la configuración inicial con \textbf{Vite} (en las tareas técnicas de infraestructura BaaS) y la configuración con \textbf{Expo + React Native Web} (para el soporte táctil y la suite de calidad).

Para evitar inconsistencias, conflictos en la integración continua (\textit{CI/CD}) y garantizar que los 6 miembros del equipo trabajen sobre el mismo entorno, se formaliza la siguiente decisión:

\begin{tcolorbox}[colback=ggray,colframe=gblue,title=\textbf{Decisión Oficial de Equipo: Adopción de Expo + React Native Web}]
    \textbf{Se adopta formalmente EXPO + REACT NATIVE WEB (\texttt{expo-router}) como stack único de la aplicación.} Queda completamente descartado el uso de Vite.
\end{tcolorbox}

\subsection{Criterios de Idoneidad para el Cliente}
\begin{itemize}[leftmargin=*]
    \item \textbf{Dispositivo objetivo prioritario:} El colegio utilizará la aplicación en \textbf{tablets en formato horizontal} por estudiantes con necesidades educativas especiales. React Native Web previene el zoom accidental, gestiona toques táctiles y optimiza el diseño en franjas.
    \item \textbf{Soporte multimedia nativo:} La aplicación requiere de avisos acústicos, síntesis de voz y renderizado de gráficos vectoriales (\textit{Time Timer} y pictogramas ARASAAC), cubiertos por \texttt{expo-audio} y \texttt{react-native-svg}.
    \item \textbf{Despliegue dual:} Se compila limpiamente para despliegue web estático en \textbf{Vercel} (\texttt{npx expo export -p web}) y permite compilar una versión instalable (APK) sin cambiar código si el centro lo solicita.
    \item \textbf{Integración Continua validada:} Los flujos de GitHub Actions (\texttt{build-and-test.yml} y \texttt{ci.yml}) están configurados para Node 22 y Jest-Expo, pasando actualmente con 100\% de éxito.
\end{itemize}

\section{Tabla Maestra del Stack Tecnológico Oficial}

Todo desarrollo nuevo en el repositorio debe respetar estrictamente esta selección tecnológica:

\begin{table}[h!]
\centering
\small
\renewcommand{\arraystretch}{1.2}
\begin{tabularx}{\textwidth}{p{3.5cm} p{4.2cm} X}
\toprule
\textbf{Ámbito} & \textbf{Tecnología Oficial} & \textbf{Detalles y Versión} \\
\midrule
Runtime / Entorno & \textbf{Node.js v22 LTS} & Fijado en \texttt{.nvmrc} y en GitHub Actions \\
Gestor de paquetes & \textbf{npm} & Mantener coherente \texttt{package-lock.json} \\
Lenguaje & \textbf{TypeScript (~6.0)} & Compilación estricta (\texttt{strict: true}) \\
Frontend Base & \textbf{Expo Web (~58.0)} & Con \texttt{react-native-web \^{}0.21} y React 19 \\
Enrutamiento & \textbf{Expo Router (~58.0)} & Enrutador basado en carpetas (\texttt{app/}) \\
Estilos e Interfaz & \textbf{StyleSheet React Native} & Contrastes accesibles WCAG 2.1 AA \\
Backend \& BaaS & \textbf{Supabase (PostgreSQL)} & Base de datos con RLS y Auth con JWT \\
Comunicación Aumentativa & \textbf{API ARASAAC} & Catálogo oficial de pictogramas \\
Testing Unitario & \textbf{Jest + Jest-Expo} & Con \texttt{@testing-library/react-native} \\
Calidad \& Accesibilidad & \textbf{ESLint 9} & Con \texttt{eslint-plugin-jsx-a11y} \\
Formateo & \textbf{Prettier} & Ejecución con \texttt{format:check} \\
CI / CD & \textbf{GitHub Actions} & Workflows obligatorios en cada Pull Request \\
Despliegue Cloud & \textbf{Vercel} & Conectado a la rama \texttt{main} \\
\bottomrule
\end{tabularx}
\caption{Matriz de tecnologías oficiales y estandarizadas.}
\end{table}

\newpage

\section{Auditoría y Revisión de Tareas Realizadas (Sprint 1)}

Se ha realizado una auditoría completa del código implementado hasta la fecha para asegurar su total alineación con el estándar Expo y los criterios de aceptación en Jira:

\subsection{HUT-01 / SCRUM-154: Infraestructura Cloud y Conexión Supabase}
\begin{itemize}[leftmargin=*]
    \item \textbf{SCRUM-155 (T-T01.1):} Despliegue continuo en Vercel. Se adaptó la configuración para permitir el empaquetado de Expo Web sin dependencias residuales de Vite.
    \item \textbf{SCRUM-157 (T-T01.3):} Conexión de cliente y variables de entorno. Se implementó retrocompatibilidad en \texttt{src/lib/supabaseClient.ts} para aceptar tanto variables con prefijo \texttt{EXPO\_PUBLIC\_} como las anteriores \texttt{VITE\_}.
\end{itemize}

\subsection{HU-02 / SCRUM-54: Inicio de Sesión Estándar para Supervisores (Golden Story)}
\begin{itemize}[leftmargin=*]
    \item \textbf{SCRUM-56 (T-02.2 Backend):}
    \begin{itemize}
        \item Servicio de autenticación JWT y contraseñas cifradas en Supabase Auth (\texttt{RNF1.3}).
        \item Control de roles RBAC (\texttt{tutor\_docente}, \texttt{familiar}, \texttt{admin}, \texttt{especialista}) con redirección dinámica (\texttt{RF1.4}).
        \item Política de seguridad contra enumeración de usuarios (\texttt{RF1.8}): mensaje genérico \textit{"Credenciales incorrectas o usuario no autorizado."}
        \item 11 tests unitarios automáticos cubriendo el 100\% de las ramas lógicas.
    \end{itemize}
    \item \textbf{SCRUM-55 (T-02.1 Frontend):}
    \begin{itemize}
        \item Pantalla de login accesible (\texttt{app/login.tsx}) desarrollada íntegramente con componentes de React Native Web.
        \item Cumplimiento WCAG 2.1 AA con etiquetas accesibles, roles semánticos y mensajes de alerta para lectores de pantalla.
        \item Validación reactiva de formulario y selector Demo para agilizar las pruebas de QA.
        \item 6 tests de integración con \texttt{userEvent} simulando interacción real.
    \end{itemize}
    \item \textbf{SCRUM-57 (T-02.3 QA):} Asignada al Gestor de Calidad (Javier Miñán) para pruebas avanzadas de expiración de sesión y estrés de tokens JWT.
\end{itemize}

\subsection{DevOps y Configuración de Calidad}
\begin{itemize}[leftmargin=*]
    \item Se homogeneizó el archivo \texttt{tsconfig.json} incluyendo \texttt{"types": ["jest"]} para que el chequeo de tipos de TypeScript (\texttt{tsc --noEmit}) sea impecable.
    \item Se actualizó \texttt{.prettierignore} para excluir artefactos de documentación y no entorpecer la validación del linter.
\end{itemize}

\section{Checklist de Entrega Obligatoria antes de cada Commit}

Para evitar que se suban Pull Requests con errores que bloqueen el pipeline de GitHub Actions, cada desarrollador \textbf{debe ejecutar localmente los siguientes 4 comandos}:

\begin{tcolorbox}[colback=ggray,colframe=gnavy,title=\textbf{Protocolo de Verificación Local (4 Checks Verdes)}]
\texttt{1. npm test -{}- -{}-ci \qquad\qquad\quad\quad (Tests unitarios con Jest)}\\
\texttt{2. npx tsc -{}-noEmit \qquad\qquad\quad\quad\ (Chequeo estricto TypeScript)}\\
\texttt{3. npm run lint \qquad\qquad\quad\quad\quad (Calidad y accesibilidad ESLint)}\\
\texttt{4. npm run format:check \qquad\quad (Formateo de código Prettier)}
\end{tcolorbox}

\noindent\textit{Nota: Si el comando de Prettier detecta desalineaciones, ejecutar simplemente \texttt{npm run format} para que corrija los archivos automáticamente antes de confirmar el commit.}

\vspace{0.8cm}
\begin{center}
\textbf{Los Gantt-sters} --- Subgrupo 2 · UGR 2026/2027\\
\textit{"Calidad, accesibilidad y rigor en el desarrollo ágil."}
\end{center}

\end{document}
